\chapter{调和函数}

\begin{solution}{harmonic-real-imaginary}
若 $f=u+iv$ 全纯，则 $u_x=v_y$、$u_y=-v_x$。再求导并用混合偏导相等，
\[
 \Delta u=u_{xx}+u_{yy}=v_{yx}-v_{xy}=0,
 \qquad
 \Delta v=-u_{yx}+u_{xy}=0.
\]
对 $f(z)=z^n$，写 $z=re^{i\theta}$，其实部为
$u=r^n\cos(n\theta)$，虚部
$v=r^n\sin(n\theta)$ 是一个调和共轭；所有调和共轭为 $v+C$。
\end{solution}

\begin{solution}{harmonic-conjugate-closed-form}
令 $\omega=-u_y\,dx+u_x\,dy$。若把它写为 $Pdx+Qdy$，则
\[
 Q_x-P_y=u_{xx}+u_{yy}=\Delta u.
\]
故 $u$ 调和时 $d\omega=0$。若 $v$ 满足 $dv=\omega$，就有
$v_x=-u_y$、$v_y=u_x$，即 Cauchy--Riemann 方程。若 $v_1,v_2$ 都满足，
则两者之差的两个偏导均为零；沿区域中连接任意两点的折线积分，差值不变，
所以它们相差常数。
\end{solution}

\begin{solution}{harmonic-log-obstruction}
对 $u(x,y)=\tfrac12\log(x^2+y^2)$，
\[
 -u_y\,dx+u_x\,dy
 =\frac{-y\,dx+x\,dy}{x^2+y^2}.
\]
沿单位圆 $x=\cos t,y=\sin t$，该形式化为 $dt$，故积分为 $2\pi$。
若存在全局调和共轭 $v$，该形式应等于 $dv$，而恰当形式沿闭曲线的积分为零，
矛盾。这也说明局部辐角不能拼成穿孔平面上的单值函数。
\end{solution}

\begin{solution}{harmonic-circle-disk-mean}
由极坐标公式和每个 $0<s<r$ 上的圆周平均，
\begin{align*}
 \int_{B(a,r)}u(x)dx
 &=\int_0^r\int_0^{2\pi}u(a+se^{it})s\,dt\,ds\\
 &=\int_0^r2\pi u(a)s\,ds
 =\pi r^2u(a).
\end{align*}
除以圆盘面积 $\pi r^2$ 即得圆盘平均。$s$ 是极坐标 Jacobian，若漏掉它便
会得到错误常数。
\end{solution}

\begin{solution}{harmonic-mean-converse}
径向核可写为 $\rho_\varepsilon(x)=R_\varepsilon(|x|)$。于是
\begin{align*}
 (u*\rho_\varepsilon)(a)
 &=\int_0^\varepsilon R_\varepsilon(r)r
 \int_0^{2\pi}u(a-re^{it})dt\,dr\\
 &=2\pi u(a)\int_0^\varepsilon R_\varepsilon(r)rdr=u(a).
\end{align*}
最后等式使用核的总积分为一。故 $u$ 局部等于光滑函数。

对 $e(t)=(\cos t,\sin t)$，Taylor 公式给
\[
 u(a+re(t))=u(a)+r\nabla u(a)\cdot e(t)
 +\frac{r^2}{2}e(t)^{\mathsf T}D^2u(a)e(t)+o(r^2),
\]
余项关于 $t$ 一致。平均时 $e$ 的一次矩为零，
$\overline{\cos^2}=\overline{\sin^2}=1/2$、
$\overline{\sin\cos}=0$，故二次项平均为 $r^2\Delta u(a)/4$。平均值恒等
迫使 $\Delta u(a)=0$。
\end{solution}

\begin{solution}{harmonic-maximum-uniqueness}
设 $u,v\in C(\overline\Omega)$ 都在 $\Omega$ 内调和且边界值同为 $g$。
$w=u-v$ 调和并在边界为零。最大值原理给
$\max_{\overline\Omega}w\le0$；对 $-w$ 使用同一结论，得
$\min_{\overline\Omega}w\ge0$。所以 $w=0$，即 $u=v$。
\end{solution}

\begin{solution}{harmonic-harnack-chain}
固定基点 $a\in\Omega$。用有限条紧含于 $\Omega$ 的折线连接 $a$ 与覆盖
$K$ 的有限个小圆盘中心。该有限并到边界有正距离，故可选统一半径 $r>0$，
使沿每条折线相邻中心距离小于 $r$，且每个中心的三倍闭圆盘仍含于
$\Omega$。设从 $a$ 到任一覆盖圆盘最多经过 $N$ 次相邻比较。正文的局部估计
每次至多引入因子 $4$，故
\[
 4^{-N}u(a)\le u(x)\le4^Nu(a)\qquad(x\in K).
\]
于是可取 $C_K=4^{2N}$ 使 $\sup_Ku\le C_K\inf_Ku$。常数不是最佳的，但其
依赖只来自有限圆盘链。
\end{solution}

\begin{solution}{harmonic-harnack-convergence}
设 $u_n$ 单调递减，并且 $(u_n(a))$ 下方有界。取调和函数
$v_n=u_1-u_n\ge0$，则 $v_n$ 单调递增，且 $v_n(a)$ 上方有界。Harnack
收敛原理给 $v_n$ 局部一致收敛到调和函数 $v$。所以
$u_n=u_1-v_n$ 局部一致收敛到调和函数 $u_1-v$。若没有一点的下界，另一种
可能是 $u_n\to-\infty$ 在紧集上一致成立。
\end{solution}

\begin{solution}{harmonic-poisson-modes}
Poisson 核的 Fourier 展开为
\[
 P_r(\theta-t)=1+2\sum_{k\ge1}r^k
 \cos k(\theta-t).
\]
利用三角函数正交性，边界数据 $\cos nt$ 与 $\sin nt$ 分别给
\[
 r^n\cos n\theta,\qquad r^n\sin n\theta.
\]
因此若
$g(t)=a_0+\sum_{n=1}^N(a_n\cos nt+b_n\sin nt)$，其调和延拓为
\[
 u(re^{i\theta})=a_0+
 \sum_{n=1}^Nr^n(a_n\cos n\theta+b_n\sin n\theta).
\]
有限和无需讨论交换极限。
\end{solution}

\begin{solution}{harmonic-poisson-boundary}
给定 $\varepsilon>0$，由 $g$ 在圆周上一致连续，取 $\delta>0$，使
$|g(e^{it})-g(e^{i\theta})|<\varepsilon$ 当圆周距离小于 $\delta$。利用
Poisson 核积分为一，
\begin{align*}
 |Pg(re^{i\theta})-g(e^{i\theta})|
 &\le\frac1{2\pi}\int P_r(\theta-t)
 |g(e^{it})-g(e^{i\theta})|dt\\
 &\le\varepsilon+
 2\|g\|_\infty\frac1{2\pi}
 \int_{|t-\theta|\ge\delta}P_r(\theta-t)dt.
\end{align*}
在后一积分区域，分母
$1-2r\cos(\theta-t)+r^2$ 有与 $r$ 接近一时无关的正下界，而分子
$1-r^2\to0$。故第二项关于 $\theta$ 一致趋零。先令 $r$ 接近一，再令
$\varepsilon\downarrow0$，得到一致边界收敛。
\end{solution}

\begin{solution}{harmonic-sharp-harnack}
先在单位圆盘处理。若 $|z|=r$、$|\zeta|=1$，则
\[
 1-r\le|\zeta-z|\le1+r,
\]
所以
\[
 \frac{1-r}{1+r}\le
 \frac{1-r^2}{|\zeta-z|^2}
 \le\frac{1+r}{1-r}.
\]
为避免预先假定 $u$ 连续到单位圆周，固定 $r<R'<1$，对
$v(\xi)=u(R'\xi)$ 使用单位圆盘的 Poisson 表示。在点 $\xi=z/R'$ 处，核的
上下界把 $r$ 换成 $r/R'$，故
\[
 \frac{R'-r}{R'+r}u(0)
 \le u(z)\le
 \frac{R'+r}{R'-r}u(0).
\]
令 $R'\uparrow1$，得到
\[
 \frac{1-r}{1+r}u(0)\le u(z)\le\frac{1+r}{1-r}u(0).
\]
平移伸缩给一般 $B(a,R)$ 的公式。

常数不可改进。对趋近某个边界点 $\zeta_0$ 的非负连续边界数据，其 Poisson
积分比值 $u(z)/u(0)$ 可任意接近核在 $\zeta_0$ 的值；让数据集中在与 $z$
同向或反向的边界点，分别逼近上、下界。
\end{solution}

\begin{solution}{harmonic-liouville}
设 $u$ 在 $\C$ 上调和且有上界。平面单连通，故 $u$ 有全局调和共轭
$v$，函数 $f=u+iv$ 整。于是
\[
 |e^{f(z)}|=e^{u(z)}
\]
有界。Liouville 定理说明 $e^f$ 为常数；求导得 $f'e^f=0$，而 $e^f$ 不为
零，所以 $f'=0$，$u$ 为常数。有下界时对 $-u$ 使用同一论证。
\end{solution}

\begin{solution}{harmonic-green-disk}
固定 $a\in\mathbb D$。在 $z\ne a$ 时，
\[
 G(z,a)=\frac1{2\pi}\bigl(\log|1-\overline a z|-\log|z-a|\bigr)
\]
是两个无零全纯函数的对数模之差，故调和。若 $|z|=1$，则
$|1-\overline a z|=|z-a|$，所以边界值为零。第一项在圆盘中光滑调和，
第二项由基本解公式给
$-\Delta[-(2\pi)^{-1}\log|z-a|]=\delta_a$，故
$-\Delta_zG=\delta_a$。

令 $z=r\zeta$、$|\zeta|=1$。由非零复函数的对数模求导公式，
\[
 2\pi\partial_rG(r\zeta,a)
 =\operatorname{Re}\left(
 -\frac{\overline a\zeta}{1-r\overline a\zeta}
 -\frac{\zeta}{r\zeta-a}\right).
\]
在 $r=1$ 时记 $b=\overline\zeta a$，括号中的两项分别为
$-\overline b/(1-\overline b)$ 和 $-1/(1-b)$。因此其实部为
\[
 -\operatorname{Re}\frac{1+b}{1-b}
 =-\frac{1-|b|^2}{|1-b|^2}
 =-\frac{1-|a|^2}{|\zeta-a|^2}.
\]
于是
\[
 -\partial_rG(r\zeta,a)|_{r=1}
 =\frac1{2\pi}\frac{1-|a|^2}{|\zeta-a|^2}.
\]
$\partial_r$ 在单位圆上就是外法向导数，所以上式为正文公式。
\end{solution}

\begin{solution}{harmonic-step-boundary}
令边界数据在上半圆为一、下半圆为零，端点值任意。Poisson 核在零点恒为一，
故
\[
 u(0)=\frac1{2\pi}\int_0^\pi1\,dt=\frac12.
\]
在非跳跃边界点，近似恒等核给径向极限等于相应常值零或一。在端点
$1$ 与 $-1$，边界数据从零跳到一；Poisson 核关于径向方向对称，两侧各占
一半质量，故径向极限为 $(0+1)/2=1/2$。这说明不连续边界数据不能期待
Poisson 解连续取得任意指定的端点值。
\end{solution}
